\section{결론}

원래 288-cell audit는 exact representability를 확립했지만 generating equation,
graph information과 head form을 교란했다. Retrospective diagnostic은 그 defect를
드러냈지만 해결하지는 못했다.

새로 동결한 prospective cohort 두 개가 이를 해결한다. Attribution cohort에서는
train-only selection이 stochastic noise 아래 two-edge graph와 head family를
회수했다. Correct graph information은 factorized, parameter-matched generic
sparse 및 더 큰 dense head 모두에서 held-out two-mechanism composition NLL을
개선했다. Seven-sparse embedding 명제는 동등한 generic support의 존재를 확립하며,
sparse procedure가 모든 world에서 그 support를 회수하여 factorized head와 정확한
동률을 만들었다. 따라서 고유한 대상은 surface notation이 아니라 typed minimal
support와 그 learned graph다.

Population 결과는 full-support primitive intervention과 \(6/7\) 미만 symmetric
corruption에서 이 motif와 typed head를 식별하며, finite bound는 modal recovery에
필요한 repeated-input coverage를 제시한다. Random-design greedy selector는 이
정리의 범위 밖이지만 120/120 회수는 post-review \(3\times3\) noise grid 전체에서
안정적이었다.

별도 stress cohort에서는 unmodeled nonadditive synergy가 모든 world에서 fitted
model을 nonexact하게 만들었지만 learned graph는 wrong routing 대비 0.18545의 NLL
effect를 유지했다. 이 결과는 exact generator--predictor equality를 graph effect의
설명에서 제거한다. Clean-room second-language implementation은 main cohort의 모든
learned fit과 analysis mean을 재현했다. 따라서 선언한 finite stochastic population
안에서 graph information, minimal structural support, composition,
misspecification과 post-export fitting/evaluation/analysis implementation
reproducibility가 각각 분리되어 식별된다.

다음 결정적 비교는 factorized head와 generic sparse head의 비교가 아니다. Entity의
생성과 삭제, identity 유지 및 relation 재구성 아래에서 동일한 dependency graph를
회수한 두 model을 비교해야 한다. 이 비교가 어떤 효과를 식별하려면 두 arm이 동일한
observable information을 받아야 하며 representational commitment만 달라야 한다.
Identity 또는 typed relation을 한 arm에만 input으로 제공하면 결과는 information
asymmetry의 산물이 되며, 이는 원래 finite-family audit에서 진단한 것과 같은 종류의
confound다. Turnover quantity \(b(p)\), 그 zero criterion과 survivor-only preservation의
vacuity 명제는 companion theory에 이미 formalize되어 있다. 따라서 turnover에 초점을
둔 연구가 이 설계의 가장 작은 self-contained instance다.
